\documentclass{article}
\bibliography{../vignettes/glmer.bib}
\usepackage{natbib}
\usepackage{bm,amsmath,amsfonts}
\newcommand{\Var}{\operatorname{Var}}
\newcommand{\abs}{\operatorname{abs}}
\newcommand{\bLt}{\ensuremath{\bm\Lambda_\theta}}
\newcommand{\mc}[1]{\ensuremath{\mathcal{#1}}}
\newcommand{\trans}{\ensuremath{^\top}}
\newcommand{\yobs}{\ensuremath{\bm y_{\mathrm{obs}}}}
\newcommand{\Lsat}{\ensuremath{L_{\mathrm{sat}}}}
\newcommand*{\eq}[1]{eqn.~\ref{#1}}% or just {(\ref{#1})}
\DeclareMathOperator*{\argmax}{arg\,max}
\DeclareMathOperator*{\argmin}{arg\,min}
\DeclareMathOperator*{\ylogy}{ylogy}

\begin{document}


\subsection{Determining the conditional mode}
\label{sec:conditionalMode}

The iteratively reweighted least squares (IRLS) algorithm is an
efficient method of determining the maximum likelihood
estimates of the coefficients in a generalized linear model.
We extend it to a \emph{penalized iteratively reweighted least squares}
(PIRLS) algorithm for determining the conditional mode, $\tilde{\bm u}_{\beta,\theta}$
(eq.~\ref{eq:condmodeGLMM}).
This algorithm has the form
\begin{enumerate}
\item Given parameter values, $\bm\beta$ and $\bm\theta$, and starting
  estimates, $\bm u_0$, of the spherical random effects, evaluate the linear predictor, $\bm\eta$, the
  corresponding conditional mean, $\bm\mu_{\mc Y|\mc U=\bm u}$, and
  the conditional variance.
  Establish the weights as the inverse of the variance.
  We write these weights in the form of a diagonal weight matrix, $\bm W$, although they are stored and manipulated as a vector.
\item Let 
  \begin{equation}
  Q(\bm u) = \left\|\bm W^{1/2}\left(\bm y-\bm\mu_{\mc
    Y|\mc U=\bm u}\right)\right\|^2+\|\bm u\|^2.
  \end{equation}

  Solve the penalized, weighted, nonlinear least squares problem
  \begin{equation}
    \label{eq:weightedNLS}
    \arg\min_{\bm u}\left( Q(\bm u) \right).
  \end{equation}
\item Update the weights, $\bm W$, and check for convergence.  If not
  converged, go to step 2.
\end{enumerate}

We use a Gauss-Newton algorithm with an orthogonality convergence
criterion~\citep[\S 2.2.3]{bateswatts88:_nonlin} to solve the
penalized, weighted, nonlinear least squares problem in step 2.
At each step of the algorithm the increment to the conditional modes
$\bm\delta^{(i)}$ 
($\bm u^{(i+1)} = \bm u^{(i)} + \bm\delta^{(i)}$)
satisfies
\begin{equation}
  \label{eq:incrEq}
  \bm P\left(\bLt\trans\bm Z\trans\bm M^{(i)}\bm W\bm M^{(i)}\bm Z\bLt+\bm I_q\right)\bm P\trans
      \bm\delta^{(i)}=\bLt\trans\bm Z\trans\bm M^{(i)}\bm W(\bm y-\bm\mu^{(i)}) - \bm u^{(i)}
\end{equation}
where $\bm u^{(i)}$ is the current value of $\bm u$, $\bm\mu^{(i)}$ is the
corresponding conditional mean of $\mc Y|\mc U=\bm u^{(i)}$ and $\bm M^{(i)}$ is
the Jacobian matrix of the vector-valued inverse link, evaluated at
$\bm\mu^{(i)}$.  That is
\begin{equation}
  \label{eq:Jacobian}
  \bm M^{(i)}=\left.\frac{d\bm\mu}{d\bm\eta\trans}\right|_{\bm\eta^{(i)}},
\end{equation}
which will be a diagonal matrix so, as for the weights, we store and
manipulate the Jacobian as a vector.

We solve (\ref{eq:incrEq}) using the sparse Cholesky factor.  At convergence, the factor, $\bm L_{\beta,\theta}$, satisfies
\begin{equation}
  \label{eq:CholFactorGLMM}
  \bm L_{\beta,\theta}\bm L_{\beta,\theta}\trans =
  \bm P\left(\bLt\trans\bm Z\trans\bm M\bm W\bm M\bm Z\bLt+\bm I_q\right)\bm P\trans,
\end{equation}
where $\bm P$ denotes a fill-reducing permutation that is computed based on the sparsity pattern
of $\bm Z \bm \Lambda_{\theta}$ prior to the optimization process (see \cite{bates2015fitting}).

When solving for the minimum of $Q(\bm u)$, the proposed update
for $\bm u$ may increase rather than decrease the value of the objective 
function $Q(\bm u)$, if the step size $\bm \delta^{(i)}$ estimated by the Gauss-Newton step
is too large. Let $Q(\bm u^{(i)})$ represent the value of the objective 
function $Q(\bm u)$ at the current iterate $\bm u_{i}$; similarly, let 
$Q(\bm u^{(i-1)}$ represent the the value of the objective function of the 
previous iteration. 

If $Q(\bm u^{(i+1)}) > Q(\bm u^{(i)})$, we proceed with a \emph{step-halving}
  procedure. That is, we successively consider scaled updates of the form
  \begin{equation}
  \bm u^{(i,j)} = \bm u^{(i)} + \frac{\bm \delta^{(i)}}{2^j}, \quad j = 1, 2, \ldots,
\end{equation}
computing $Q(\bm u^{(i,j)})$ at each step. The step-halving procedure 
continues until a value $j^*$ of $j$ is found such that 
$Q(\bm u^{(i,j^*)}) < Q(\bm u^{(i)})$,
after which the iterate is then updated by setting $\bm u^{(i+1)} = \bm u^{(i,j^*)}$.

If no such $j$ is found after a predetermined number of halvings, then the 
step-halving procedure has failed and the algorithm terminates. 

As we show in the next section, the matrix
$(\bm L_{\beta,\theta}\bm L_{\beta,\theta}\trans)^{-1}$ is a Laplace approximation of the covariance matrix for the spherical random effects, conditional on the
observed data. This fact is useful for constructing a nonlinear
objective function for finding the approximate maximum likelihood
estimates of $\theta$ and $\beta$.


\end{document}
